\documentclass[12pt]{article}
\usepackage{amsthm,amssymb,amsmath,amsfonts}
\usepackage[a4paper, top=25mm, bottom=30mm, left=25mm, right=25mm]{geometry}
\usepackage[pagebackref=false,colorlinks,linkcolor=black,citecolor=black]{hyperref}
\usepackage[nameinlink]{cleveref}
 \AtBeginDocument{%
    \crefname{equation}{برابری}{equations}%
    \crefname{chapter}{فصل}{chapters}%
    \crefname{section}{بخش}{sections}%
    \crefname{appendix}{پیوست}{appendices}%
    \crefname{enumi}{مورد}{items}%
    \crefname{footnote}{زیرنویس}{footnotes}%
    \crefname{figure}{شکل}{figures}%
    \crefname{table}{جدول}{tables}%
    \crefname{theorem}{قضیه}{theorems}%
    \crefname{lemma}{لم}{lemmas}%
    \crefname{corollary}{نتیجه}{corollaries}%
    \crefname{proposition}{گزاره}{propositions}%
    \crefname{definition}{تعریف}{definitions}%
    \crefname{result}{نتیجه}{results}%
    \crefname{example}{مثال}{examples}%
    \crefname{remark}{نکته}{remarks}%
    \crefname{note}{یادداشت}{notes}%
    \crefname{observation}{مشاهده}{observations}%
    \crefname{algorithm}{الگوریتم}{algorithms}%
    \crefname{cproof}{برهان}{cproofs}%
}

\usepackage{tikz}
\usepackage{graphicx}
\usepackage{color}

\usepackage{setspace}
\doublespacing
\usepackage{listings}

\lstnewenvironment{code}[1][]%
  {\minipage{\linewidth} 
   \lstset{basicstyle=\ttfamily\footnotesize,frame=single,#1}}
  {\endminipage}
  
\usepackage{titletoc}
\usepackage{tocloft}
\usepackage{enumitem}
\usepackage{lipsum} % for dummy text
\usepackage{tikz}
\usepackage{eso-pic}

\AddToShipoutPictureBG{% Add picture to background of every page
    \begin{tikzpicture}[remember picture,overlay]
    \draw[line width=1pt]
        ([shift={(1cm,-1cm)}]current page.north west)
        rectangle
        ([shift={(-1cm,1cm)}]current page.south east);
    \end{tikzpicture}
}



\usepackage{algorithm}
% \usepackage[noend]{algpseudocode}
\usepackage[noend]{algorithmic}
\renewcommand{\algorithmicrequire}{\textbf{Input:}}
\renewcommand{\algorithmicensure}{\textbf{Output:}}

\usepackage{tabularx}
\makeatletter
\newcommand{\multiline}[1]{%
  \begin{tabularx}{\dimexpr\linewidth-\ALG@thistlm}[t]{@{}X@{}}
    #1
  \end{tabularx}
}
\makeatother

\usepackage{float}
\usepackage{verbatim}
\makeindex
\usepackage{sectsty}
\usepackage{xepersian}
\SepMark{-}
\settextfont[Scale=1.2,Path=fonts/,BoldFont=B Nazanin Bold.ttf]{B Nazanin.ttf}
\setlatintextfont{Times New Roman}
\renewcommand{\labelitemi}{$\bullet$}

\theoremstyle{definition}
\newtheorem{definition}{تعریف}[section]
\newtheorem{remark}[definition]{نکته}
\newtheorem{note}[definition]{یادداشت}
\newtheorem{example}[definition]{نمونه}
\newtheorem{question}[definition]{سوال}
\newtheorem{remember}[definition]{یاداوری}
\newtheorem{observation}[definition]{مشاهده}
\theoremstyle{theorem}
\newtheorem{theorem}[definition]{قضیه}
\newtheorem{lemma}[definition]{لم}
\newtheorem{proposition}[definition]{گزاره}
\newtheorem{corollary}[definition]{نتیجه}
\newtheorem*{cproof}{برهان}



\begin{document}


\fontsize{12pt}{14pt}\selectfont

\begin{minipage}{0.1\textwidth}
\includegraphics[width=2cm]{etc/aut}
\end{minipage}%
\hfill%
\begin{minipage}{0.6\textwidth}\centering
\fontsize{10pt}{10pt}\selectfont
به‌نام خدا \\
پوشش دیسک واحد \\
نام درس \\
نام و نام خانوادگی \\
\vspace{0.25cm}
\begingroup
\fontsize{8pt}{8pt}\selectfont
دانشگاه صنعتی امیرکبیر، دانشکده مهندسی کامپیوتر \\
خردادماه 1401 \\
\endgroup
\end{minipage}%
\hfill%
\begin{minipage}{0.1\textwidth}
\includegraphics[width=1.8cm]{etc/ce}
\end{minipage}

\vspace{0.5cm}

\noindent\rule{\textwidth}{1pt}

\begin{abstract}
\noindent
این پروژه برای الانیلی عددی
\end{abstract}

\section{شرح مسئله}

معادله غیر خطی تیر اویلر-برنولی با مقطع یکنواخت، در حالت استاتیکی را بصورت زیر در نظر بگیرید:
\begin{equation}
	E I \frac{\partial^4 w}{\partial x^4}-\frac{3}{2} E A\left(\frac{\partial w}{\partial x}\right)^2\left(\frac{\partial^2 w}{\partial x^2}\right)=q(x)
\end{equation}
که در آن$E$ مدول الاستیسیته، $I$ ممان اینرسی مقطع تیر، $w(x)$ خیز تیر،$q(x)$ بار گسترده، $A$ مساحت مقطع و $L$ طول تیر میباشد.
\\
\begin{itemize}
  \item هندسه تیر: مقطع تیر مستطیلی به ابعاد $0.1m * 0.2m$ (ارتفاع*عرض) و طول تیر $2m$ 
  \item جنس تیر: آلومینیوم، $E = 70GPa$
  \item نیرو ها: $q(x)=10^4 * a\left(\frac{x}{L}\right)^b N / m$
\end{itemize}
از آن جایی که شماره دانشجویی من ۴۰۰۱۲۶۰۳۴  هست پس طبق فرض مسيله تکیه گاه های مسيله من یک سر گیر دار و یک سر تکیه گاه ساده در نظر گرفته می‌شود.
\begin{itemize}
	\item $a=3$
	\item $b=4$
\end{itemize}
	
\section{ روش حل:}
\subsection{معرفی روش های حل:}
\subsubsection{روش \lr{Method of Adjoints}}

%code space
\begin{latin}
\begin{verbatim}
print("Hello, world!")
the new command 
\end{verbatim}
\end{latin}

\subsubsection{روش \lr{GDQ}}

%code space
\begin{latin}
\begin{verbatim}
print("Hello, world!")
the new command 
\end{verbatim}
\end{latin}

\end{document}